% !PRINTED 67
\weberpage{67}
Dass im Falle eines geraden $\beta$ jeder der beiden Factoren verschwindet, sieht man, wenn man in (15)
\[
n' = 1 + 2^{\lambda-1},
\]
also
\[
r^{n'} = -r, \qquad \nu_1' = 0, \qquad \nu' = 2^{\lambda-3}
\]
setzt. Dann giebt diese Formel, da $n$ ungerade ist:
\[
\bigl((-1)^{\beta_1}, \Theta^\beta, r\bigr) = -\Theta^{\beta 2^{\lambda-3}} \bigl((-1)^{\beta_1}, \Theta^\beta, r\bigr).
\]
Bei geradem $\beta$ ist aber $\Theta^{\beta 2^{\lambda-3}} = +1$ und folglich:
\weq{18}{\bigl((-1)^{\beta_1}, \Theta^\beta, r\bigr) = 0;}
also:

2.~\gesperrt{Ist der Grad der Einheitswurzel $r$ eine Potenz von 2 und grösser als 4, so verschwindet die Resolvente $\bigl((-1)^{\beta_1}, \Theta^\beta, r\bigr)$ dann und nur dann, wenn $\beta$ gerade ist.}

Im Falle $m=4$, also $r=i$, hat man nur die zwei Resolventen $i+i^3, i-i^3$, die wir unter der Bezeichnung $((-1)^\beta, i)$, $\beta=0,1$ zusammenfassen können, und es ist $((-1)^\beta, i)$ dann und nur dann $=0$, wenn $\beta=0$ ist.

\begin{center}
\S.\,18.\\
\gesperrt{Kreistheilungsk\"orper.}
\end{center}

Die Theorie der \gesperrt{Abel'schen} Gruppen eröffnet uns einen tieferen Einblick in die Theorie der Einheitswurzeln und der daraus entspringenden algebraischen Zahlen.

Es möge jetzt $m$ irgend eine ganze positive Zahl sein, die in ihre Primfactoren zerlegt sei:
\weq{1}{m = 2^\lambda q_1^{\varkappa_1} q_2^{\varkappa_2} \cdots,}
und es sei $r$ eine primitive $m^{\text{te}}$ Einheitswurzel. Den Fall $\lambda=1$ können wir ein- für allemal von unserer Betrachtung ausschliessen; denn wenn $m$ ungerade ist und $r$ die primitiven $m^{\text{ten}}$ Einheitswurzeln durchläuft, so kommen darunter keine zwei entgegengesetzte vor, und $-r$ durchläuft die primitiven $2m^{\text{ten}}$ Einheitswurzeln.

$r$ ist die Wurzel einer ganzzahligen irreduciblen \gesperrt{Abel'schen} Gleichung vom Grade
\weq{2}{\nu = \varphi(m),}
wie wir im \S.\,134 des ersten Bandes nachgewiesen haben.
